\par
\SourceLeaf{366}%
\pdfbookmark[0]{Anzeigen. Herschel: Cotesischer Lehrsatz. 1814.}{anzeigen-herschel-1814}
\begin{center}
\rule{.85\textwidth}{.3pt}\par
Göttingische gelehrte Anzeigen.\quad 1814 Februar 14.\par
\rule{.85\textwidth}{.3pt}
\end{center}
\hspace*{\parindent}\textit{Philosophical Transactions of the Royal Society of London for the Year} 1813. IV und 304 Seiten. 26 S. \textit{Meteorological Journal} und 8 S. \textit{Index} in Quart.

\textit{Mathematische und astronomische Abhandlungen. --- Über eine merkwürdige Anwendung des} \textsc{Cotes}\textit{ischen Lehrsatzes, von} \textsc{J. F. W. Herschel} (Sohn des Astronomen). Es sei \(n\) eine beliebige ganze Zahl, \(n\omega=360^0\) und \(N\) irgend ein Winkel. Unter diesen Voraussetzungen gibt der \textsc{Cotes}ische Lehrsatz das Product aus allen Radiis Vectoribus, denen in einem Kegelschnitt, nach astronomischer Art zu reden, die wahren Anomalien \(N,\ N+\omega,\ N+2\omega,\ N+3\omega\mathpunct{.}\mathpunct{.}\mathpunct{.}\mathpunct{.}N+(n-1)\omega\) entsprechen, durch einen einfachen Ausdruck. Wenn gleich diese und andere ähnliche Entwickelungen, welche den Gegenstand des Aufsatzes ausmachen, an sich keine besondere Schwierigkeiten haben, so liest man diesen doch mit Vergnügen wegen der Art der Behandlung. Was der Verf. über die Bezeichnung \(\cos^2 A\) sagt, welches einige neuere mathematische Schriftsteller für das Quadrat von \(\cos A\), ganz gegen alle Analogie, gebrauchen, da es dieser zufolge den Cosinus eines Bogens \(=\cos A\) bedeuten sollte, hat ganz unsern Beifall.

\pdfbookmark[0]{Anzeigen. Mollweide: Commentationes mathematico-philologicae. 1814.}{anzeigen-mollweide-1814}
\begin{center}
\rule{.85\textwidth}{.3pt}\par
Göttingische gelehrte Anzeigen.\quad 1814 Mai 2.\par
\rule{.85\textwidth}{.3pt}
\end{center}
\hspace*{\parindent}\textit{Commentationes mathematico-philologicae tres, sistentes explicationem duorum locorum difficilium, alterius} \textsc{Virgilii}, \textit{alterius} \textsc{Platonis}, \textit{itemque examinationem duorum mensurarum praeceptorum} \textsc{Columellae}. \textit{Adjecta est epistola ad v.\,cl.} \textsc{J. G. Schneider} \textit{de excerptis geometricis} \textsc{Epaphroditi} \textit{et} \textsc{Vitruvii Rufi} \textit{scripta ab auctore harum commentationum} \textsc{Carolo Brandano Mollweide}, \textit{astron. in acad. Lipsiensi professore.} Leipzig 1813. 122 Seiten in Octav, nebst einer Kupfertafel.

[Die Anzeige der ersten Abhandlung ist in dem Bande für Astronomie der Werke von \textsc{Gauss} abgedruckt.]

Die zweite Abhandlung, über eine dunkle Stelle in \textsc{Plato}'s Menon, war schon im Jahre 1805 der hiesigen Königl. Soc. der Wissenschaften handschrift%
